%% =====================================================================
%%  Bounded by Specification
%%  8-page draft for IEEE RA-L (also fits ICRA/IROS).
%%
%%  CLAIM: a skill specification pairing LTL invariants with per-phase
%%  deadlines and progress budgets, calibrated from nominal runs, gives a
%%  controller-independent monitor with a verdict horizon fixed by the
%%  specification and a fault that names the stage and guard.
%%
%%  The longer argument (15 pages) is kept in ral_extended.tex.
%%  Unmeasured quantities are \TODO.  See README.md and DEFENCE.md.
%% =====================================================================

\documentclass[letterpaper, 10pt, journal, twoside]{IEEEtran}

\usepackage[T1]{fontenc}
\usepackage{amsmath,amssymb,amsfonts}
\usepackage{siunitx}
\usepackage{graphicx}
\usepackage{booktabs}
\usepackage{url}
\usepackage[hidelinks]{hyperref}
\usepackage{xcolor}

\newcommand{\TODO}[1]{\textcolor{red}{\textbf{[TODO: #1]}}}

\newtheorem{proposition}{Proposition}
\newcommand{\proofpar}{\par\noindent\emph{Proof.}\ }
\newcommand{\qedbox}{\hfill$\blacksquare$\par\medskip}

\newcommand{\G}{\mathbf{G}}
\newcommand{\F}{\mathbf{F}}
\newcommand{\X}{\mathbf{X}}

\title{Bounded by Specification: Calibrated Per-Phase Deadlines for
Independent Runtime Monitoring of Robot Skills}

\author{\TODO{authors, affiliations, funding, RA-L footnotes}}

\begin{document}
\maketitle

\begin{abstract}
A runtime monitor for a robot skill is usually handed the skill's goal
as a temporal-logic formula. That formula is a guarantee property: it
has no bad prefix, so a monitor can confirm it and never refute it, and
a stalled robot leaves it inconclusive forever. In practice the gap is
filled by the controller timing itself out, which fails exactly when the
controller is wrong about its own state. We present a skill
specification that pairs Linear Temporal Logic (LTL) invariants with a
ladder of \emph{phases}, each carrying a precondition, an invariant, a
progress condition with a violation budget, an exit condition and a
deadline. The monitor that runs it is independent of the controller it
judges, reaches a verdict within a horizon fixed by the specification,
and names the stage and guard that failed. The deadlines are not
hand-tuned: they are calibrated as conformal quantiles of phase
durations in successful runs, so the false-alarm level is stated rather
than hoped for. We evaluate the monitor as an intervention rule ---
detection delay at a fixed false-alarm level, break-even false-alarm
rate and execution saved --- on teleoperated humanoid pick-and-place and
on humanoid navigation, against a flat specification, a mission-level
deadline and per-stage bounded LTL. \TODO{headline result}
\end{abstract}

\begin{IEEEkeywords}
Formal methods in robotics and automation, failure detection and
recovery, runtime verification, humanoid robots.
\end{IEEEkeywords}

% =====================================================================
\section{Introduction}
% =====================================================================

\IEEEPARstart{A}{humanoid} executing a waypoint mission wedges itself
against a lip its traversability estimate called flat. The legs keep
cycling, the planner keeps reporting that it is following a path, and
the goal is never reached. What should a runtime monitor say?

The skill's requirement, in LTL, is something like
\begin{equation}
\label{eq:flat}
\F\big(\mathit{armed} \wedge \F(\mathit{planning} \wedge
  \F(\mathit{walking} \wedge \F \mathit{finished}))\big),
\end{equation}
which is the goal formula our own platform shipped with. A monitor for
\eqref{eq:flat} reports \textsc{inconclusive} on this trace, and on
every later tick. It is not late; it is correct. The formula is a
guarantee property: at any finite point $\mathit{finished}$ could still
occur, so there is no bad prefix and nothing to
refute~\cite{manna1990hierarchy,bauer2011runtime}. The formula that says
what a skill is \emph{for} is the one that can never report that the
skill failed.

Practice fills the gap with timeouts, and they live in the controller: a
navigation stack's progress checker, a behaviour tree's timeout
decorator, an action server's deadline. These work until the controller
is wrong about itself, which is the case above --- any check that reads
the planner's status sees \texttt{following}.

We move the bound out of the controller and into the specification. A
skill is described as cross-cutting LTL invariants plus a ladder of
\textbf{phases}; each phase declares what must hold, what counts as
progress and how long it may take. Figure~\ref{fig:two} shows one
mission both ways. The resulting monitor reads the robot rather than the
planner, cannot wait forever, and reports
\texttt{phase:tracking:timeout} rather than a failed conjunction.

A bound must come from somewhere, and a hand-picked one invites the
charge that it was tuned to the failures. We calibrate each phase's
deadline from successful runs as a conformal quantile of its duration,
which fixes the false-alarm level in advance. And since an alarm is only
worth what the response to it achieves~\cite{ahmedt2026rescue}, we
evaluate the monitor as an intervention rule, not as a classifier.

\textbf{Contributions.}
\begin{enumerate}
  \item \emph{Analysis.} A skill's goal formula is a guarantee property
  and is never refuted online under LTL$_3$; neither a finite-trace nor
  a four-valued semantics repairs this. A bound per stage does, with a
  proved horizon on time to verdict (Section~\ref{sec:spec}).
  \item \emph{Design.} A two-layer specification --- LTL automata for
  cross-cutting invariants, a phase ladder for bounded, localised and
  categorised progress obligations --- whose bounds are calibrated from
  nominal runs at a stated false-alarm level.
  \item \emph{Independence.} The monitor is an artifact separate from
  the controller and grounded in the robot's own signals; we exhibit
  the case in which the controller's own progress check is silent.
  \item \emph{Evaluation as an intervention rule} on two skill families
  on one humanoid, against three baselines. \TODO{result}
\end{enumerate}

\begin{figure}[t]
\centering
\small
\begin{tabular}{@{}p{0.45\columnwidth}p{0.45\columnwidth}@{}}
\toprule
\textbf{Flat LTL} & \textbf{Phase ladder} \\
\midrule
$\F(\mathit{armed} \wedge \F(\mathit{planning}$
$\wedge\, \F(\mathit{walking} \wedge \F \mathit{finished})))$
\vspace{3pt}

Stalled trace: \textsc{inconclusive}, forever. No diagnosis.
&
\textsc{planning}: deadline 30, budget 5\\
\textsc{tracking}: deadline 120, budget 5\\
\textsc{confirm}: deadline 10
\vspace{3pt}

Stalled trace: \texttt{phase:tracking:timeout} $\rightarrow$
\textsc{replan}.
\\
\bottomrule
\end{tabular}
\caption{One navigation mission, two encodings, from the same
specification file (bounds in ticks). \TODO{replace with a drawn figure
and an annotated stall trace.}}
\label{fig:two}
\end{figure}

% =====================================================================
\section{Background: LTL, and Why LTL}
\label{sec:background}
% =====================================================================

\textbf{LTL in brief.} LTL~\cite{pnueli1977temporal} extends
propositional logic over atomic propositions $AP$ with operators over
discrete time: $\X\varphi$ (next), $\G\varphi$ (always), $\F\varphi$
(eventually) and until. A model is an infinite word over $2^{AP}$; here
one position is one monitor tick. A formula is compiled into an
automaton whose state the monitor tracks as observations
arrive~\cite{vardi1986automata}; we use Spot~\cite{duretlutz2016spot}.
Two property classes matter~\cite{manna1990hierarchy}. A \emph{safety}
property, such as $\G\neg\mathit{collision\_risk}$, is violated by a
finite bad prefix. A \emph{guarantee} property, such as
\eqref{eq:flat}, has none: every finite trace has an extension that
satisfies it.

\textbf{Monitoring a finite observation.} LTL$_3$~\cite{bauer2011runtime}
maps a finite prefix to $\top$ if every infinite continuation satisfies
the formula, $\bot$ if none does, and $?$ otherwise. It is the semantics
of observing an open-ended run, which is the monitor's situation at
every tick, and we adopt it for the verdict that matters: our
\textsc{violated} is LTL$_3$'s $\bot$, reported when the automaton
enters a non-accepting sink, and it is permanent. (Our
\textsc{accepted} status only reports an accepting current state and is
not LTL$_3$'s $\top$; nothing acts on it.) Two alternatives do not help
online. LTL$_f$~\cite{degiacomo2013linear} would declare
$\F\,\mathit{finished}$ false on a finite trace, but only at the last
position, and it cannot say when the trace ends --- something outside
the logic must, and in our system that is a deadline.
RV-LTL~\cite{bauer2010comparing} refines $?$ into ``presumably true''
and ``presumably false'', but $\F\,\mathit{finished}$ is presumably
false on every prefix of a \emph{successful} run until its last tick, so
the value does not separate stalled from healthy and is not monotone.

\textbf{Why LTL as the description language.} We need three things at
once. \emph{A declared property rather than an executable.} Behaviour
trees and state machines are how robot behaviour is structured, and
runtime verification has been built around
them~\cite{colledanchise2018behavior,colledanchise2021formalizing,ahmad2025unified};
but the property a tree enforces is implicit in its traversal, and in
practice the tree is the controller, so a monitor derived from it shares
the controller's view of the world. \emph{A verdict that cannot be
retracted.} A violated safety formula sits in a sink state, so a
supervisor can latch on it; a learned or language-model judgement can
change on the next frame. \emph{A property class one can reason about},
which is what exposes the unrefutable goal in the first place. Signal
Temporal Logic~\cite{maler2004monitoring,fainekos2009robustness} is the
serious alternative and does offer bounded operators; we keep the
temporal layer Boolean because a robustness value says how badly rather
than which obligation, and we push thresholds into proposition
grounding. Nothing in the phase structure precludes STL guards.

% =====================================================================
\section{Related Work}
\label{sec:related}
% =====================================================================

\textbf{Stage structure in monitoring.} That execution monitoring needs
a notion of stage is old~\cite{pettersson2005execution,degiacomo1998execution}
and now common. It is \emph{learned} from labelled episodes: RoboMonitor
predicts phase, failure and completion with a vision--language model and
must suppress spurious phase switching~\cite{ajith2026robomonitor};
ConditionNET learns preconditions and effects and compares expected with
predicted phase~\cite{sliwowski2025conditionnet}. It is \emph{consumed}:
AHA and Guardian verify a given sub-task from
keyframes~\cite{duan2025aha,pacaud2025guardian}, and REFLECT explains a
failure afterwards~\cite{liu2023reflect}. It is \emph{generated}: DoReMi
has a language model emit per-skill constraints that a vision model
checks during execution~\cite{guo2024doremi}, and Code-as-Monitor
compiles during-execution and on-completion constraints into
code~\cite{zhou2024codeasmonitor}, a split close to our invariant and
exit condition. And it is used to make \emph{planning} tractable through
hierarchical LTL~\cite{xu2024nl2hltl2plan}, whose sub-formulas remain
LTL and so inherit the property class of their parent. In the systems we
examined the per-stage obligations are states to maintain; where a time
bound appears it is a single task-level timeout.

\textbf{Controller-side progress checks.} Horizons on progress do exist,
inside controllers. Nav2's progress checker fails a controller that has
not moved a set distance within a set
time~\cite{macenski2020marathon}; behaviour trees have timeout
decorators~\cite{colledanchise2018behavior}; a steps-since-progress
counter is a standard monitoring
baseline~\cite{ahmedt2026rescue}. We do not claim the mechanism. We
claim its placement --- outside the component it judges, declared per
stage in a specification, producing a categorised verdict --- and its
calibration. The risk of the alternative is documented outside robotics
too: language-model agents asked to report their own progress are
unreliable mid-task, and such reports should be checked against
independent state~\cite{wang2026progressbar}.

\textbf{Learned failure detection and its evaluation.} A large
literature detects failures of learned policies from their own
statistics or from visual
judgement~\cite{agia2024unpacking,periasami2026frames}. Two results from
it shape our evaluation. FAIL-Detect calibrates its threshold by
conformal prediction on successful runs only, obtaining a false-positive
guarantee without failure data~\cite{xu2025faildetect}; we apply the
same idea to phase deadlines. RESCUE shows that detection and timing
metrics admit monitors that are harmful once false alarms and
intervention cost are counted, and gives a break-even false-alarm rate
computable from execution logs~\cite{ahmedt2026rescue}; we report it.
Granularity matters in neighbouring fields as well: stage-wise
evaluation of long-horizon agents discriminates better than
whole-trajectory scoring~\cite{shi2026dynsteer}.

\textbf{Specifications from language.} Translating instructions into
temporal logic is well
studied~\cite{liu2023grounding,chen2023nl2tl,fuggitti2023nl2ltl,sundarsingh2025conformalnl2ltl}.
Our specifications can be generated from a skill description and checked
against a declared sensor schema; we do not evaluate that here.

\textbf{Position.} Stage structure is prior art, and so is the timeout.
What we add is a specification in which each stage carries a calibrated
bound, a monitor for it that is independent of the controller, the
property-class argument for why the bound is necessary rather than
convenient, and an evaluation of the result as an intervention rule.

% =====================================================================
\section{A Phase-Structured Specification}
\label{sec:spec}
% =====================================================================

\subsection{The gap}

A skill specification contains two kinds of formula. Cross-cutting
invariants such as $\G\neg\mathit{collision\_risk}$ and
$\G\,\mathit{upright}$ are safety properties and are refutable. The goal
formula \eqref{eq:flat} is a guarantee property, so its LTL$_3$ verdict
is $\top$ or $?$, never $\bot$. A robot that is wedged, circling or
standing still produces an unbounded run of inconclusive verdicts. A
\emph{terminal failure condition} does not close the gap if it is
written over the controller's status, as ours was.

\subsection{Phases}

A phase is a named stage with five guards and two budgets: an
\emph{enter condition}; a \emph{precondition}, checked once on entry; an
\emph{invariant}, checked every tick; a \emph{progress condition} with a
\emph{violation budget} $m$, exhausted by $m$ consecutive failures; an
\emph{exit condition}, which advances the ladder; and a \emph{deadline}
$k$ in ticks. Each breach is a named fault,
\texttt{phase:$\langle$name$\rangle$:$\langle$kind$\rangle$}, with a
category.

Bounding an obligation is what changes its class: ``this phase is
eventually left'' has no bad prefix, while ``this phase is left within
$k$ ticks'' is decided in $k+1$ positions. That is a fact about the
bound. For the ladder it yields a horizon on time to verdict.

\begin{proposition}
\label{prop:bounded}
Let a ladder have phases $1 \dots n$ visited in order, each declaring a
finite deadline $k_i$. Within $\sum_{i=1}^{n} (k_i + 1)$ ticks of
entering the first phase, the monitor has completed the ladder or raised
a fault.
\end{proposition}

\proofpar The ladder occupies one phase at a time. A phase's tick
counter is reset on entry and incremented on every tick the phase
remains active. On any tick at which it has reached $k_i$ the machine
raises a terminal timeout fault before consulting the progress or exit
condition. So a phase is occupied for at most $k_i + 1$ ticks; phases
are left only forwards and never re-entered. \qedbox

The proposition is elementary, and we do not offer it as a result: it
says only that non-completion always yields a verdict, not that the
verdict is timely or deserved. Its value is what it rests on. The
timeout path counts ticks and evaluates no expression, so it holds
whatever the sensors or the controller report. The progress budget is
the opposite: it is what brings typical delay far below the bound, and
it is only as independent as the proposition it reads. Whether the
verdict is timely, and what it costs, is empirical
(Section~\ref{sec:eval}).

What the phases add beyond the bound: a progress budget needs a stage to
be scoped to; an overrun is caught in the stage where it occurs; and the
stage index answers the operator's first question and gives one place to
grade a fault into an intervention. Per-stage \emph{bounded} LTL,
$\G(\mathit{armed} \wedge \neg\mathit{planning} \rightarrow
\F_{\leq k}\,\mathit{planning})$, expresses the deadlines without a
phase machine, and we do not claim otherwise; it gives no stage index,
its automata grow with $k$, and scoping a budget still needs a stage
indicator. We treat it as a baseline, not an argument.

\subsection{Calibrating the bounds}
\label{sec:calib}

A deadline's false-alarm rate is the probability that a healthy run
overstays it, so choosing $k_i$ is choosing a quantile of the phase's
nominal duration. We make that explicit. From $N$ successful runs we
record each phase's duration $d_i^{(1)}, \dots, d_i^{(N)}$ and set
\[
k_i \;=\; \text{the } \lceil (N+1)(1-\alpha_i) \rceil\text{-th smallest
of } d_i^{(1)}, \dots, d_i^{(N)},
\]
the split-conformal quantile used for failure-detection thresholds
in~\cite{xu2025faildetect}. If a new healthy run is exchangeable with
the calibration runs, the phase overstays $k_i$ with probability at most
$\alpha_i$; taking $\alpha_i = \alpha / n$ bounds the episode-level
false-alarm probability from deadlines by $\alpha$. The progress budget
$m$ is calibrated the same way from the longest run of consecutive
progress violations in each successful run. Two caveats are intrinsic.
The guarantee is marginal and needs $N \geq (1-\alpha_i)/\alpha_i$ runs
to be non-trivial --- below that $k_i$ is simply the longest duration
seen. And exchangeability fails when the operator, the scene or the
robot changes, which is why the level is checked on held-out runs rather
than assumed.

\subsection{What it does not fix}

The bounds and grounding thresholds are calibration, not logic: a wrong
one yields a verified wrong answer. The ladder is a total order and does
not express concurrent stages. It cannot flicker between stages, but a
one-tick glitch in an exit condition advances it irreversibly. Finally,
a proposition that cannot be evaluated is handled inconsistently today:
a phase guard falls back to a pass, while the automata read a missing
proposition as false, so a safety fault can be hidden or fabricated
depending on how its proposition is named. A three-valued treatment that
freezes the monitor and raises a data-loss fault is designed and not
deployed.

% =====================================================================
\section{The Monitor}
\label{sec:system}
% =====================================================================

\textbf{Artifacts.} A \emph{sensor schema} declares the keys a
specification may mention. An \emph{adapter descriptor}, one per
embodiment, maps topics and fields to those keys and declares how
samples within a tick fold to one value. A \emph{skill specification}
holds the propositions with their grounding rules, the LTL invariants
with fault categories, and the phase ladder. The monitor reads a
specification and an observation stream and never the descriptor, so
nothing in it names a robot; the same engine runs both skills in
Section~\ref{sec:eval}.

\textbf{One tick.} The monitor is driven by a clock, not by message
arrival, so a bound in ticks means the same thing live and on replay.
On each pulse: (i)~each key's samples fold to one value; (ii)~the
propositions the monitor currently needs are grounded --- a rule of the
form ``True when $\langle$expression$\rangle$'' is evaluated directly,
and a proposition that is not reducible to a threshold is sent to a
model on a worker thread; (iii)~each automaton takes one step; (iv)~the
active phase's guards are consulted in the order invariant, deadline,
progress, exit; (v)~one verdict is published with the phase, each
formula's status, the truth of each guard, the violation counters and
any fault.

\textbf{From fault to intervention.} A monitor alone changes no outcome.
A fault carries a category: timeout and progress are fixed, and each
phase's invariant and precondition declare theirs. Category, imminence
and confidence select a rung of a six-step ladder --- continue, warn,
slow, replan, halt, abort --- graded in one place. A confirmed safety
fault halts or aborts; a confirmed timeout or stall requests a replan; a
fault resting on low-confidence sensing is reduced to a warning. A
supervisor overrides the base with zero velocity while a halting fault
is active and does nothing until enabled, which provides a
detection-only arm.

% =====================================================================
\section{Evaluation}
\label{sec:eval}
% =====================================================================

\subsection{Data}

\textbf{Teleoperated pick-and-place (main).} \TODO{$N$} episodes of a
Unitree G1 picking and placing a bottle under teleoperation, with
egocentric video and joint values per frame, labelled \emph{successful},
\emph{failed} or \emph{almost failed}. The skill is naturally staged ---
reach, grasp, lift, transport, place --- and is the task family most of
the related work uses. Propositions are grounded from joint values
through forward kinematics (hand position, gripper closure, hand
activity) and, for object state, from video. Phase boundaries in
successful episodes and the onset frame in failed ones are annotated
from video alone, so no label is derived from a signal the monitor
reads.

\textbf{Humanoid navigation (supporting).} Waypoint missions on the same
robot and in MuJoCo~\cite{todorov2012mujoco} simulation, with stalls
induced at known times. This is the setting in which a controller-side
check exists and can be silent.

\textbf{Why these and not a public benchmark.} We know of no public set
with multi-stage skills, stage annotations and controller-status
signals together. Manipulation failure sets built by perturbing
demonstrations~\cite{duan2025aha,pacaud2025guardian,liu2023reflect} lack
the last; the closest navigation set, ForEnt, records quadruped
entrapments in single-stage traversal~\cite{kirdwichai2026forent} and is
a natural next test of the progress budget. The recordings here have a
limit we state now: nothing in them can be intervened on, so we measure
what a monitor would have saved, not task success.

\subsection{Arms and metrics}

All arms replay identical observation streams and share propositions and
invariants. \textbf{B1}: flat goal formula. \textbf{B2}: B1 plus one
mission deadline, calibrated like ours. \textbf{B3}: per-stage bounded
LTL without a phase machine. \textbf{C}: the phase ladder. For
navigation, \textbf{B4}: the controller's own progress check.

We report, per failure class:
\emph{false-alarm rate} on held-out successful episodes, against the
level $\alpha$ fixed before any failed episode was examined;
\emph{detection delay} from annotated onset, as median and worst case
over detected episodes, with the detection rate given separately because
delay is censored when nothing is detected;
\emph{localisation}, whether the fault names the annotated stage;
\emph{execution saved}, frames from fault to episode end; and the
\emph{break-even false-alarm rate} of~\cite{ahmedt2026rescue}, the level
above which acting on the monitor costs more than ignoring it.
For almost-failed episodes we report the highest intervention rung
reached: the intended behaviour is a warning that clears, not a terminal
fault. We do not report AUROC, which scores frames and is blind to when
an alarm arrives.

Detection rate is not a finding for B1 or C. B1 detects no progress
failure by the argument of Section~\ref{sec:spec}, and C raises a fault
on every non-completing episode by Proposition~\ref{prop:bounded}. Both
rates are reported as a check on the implementation. The empirical
questions are whether the calibrated false-alarm level holds, how far
below the bound the delay falls, and whether C improves on B2 and B3.

\subsection{Results}

\begin{table}[t]
\centering
\caption{Pick-and-place, $\alpha = \TODO{}$. Delay in frames.}
\label{tab:main}
\small
\begin{tabular}{@{}lccccc@{}}
\toprule
arm & false & near-miss & det. & median & break- \\
    & alarm & fault     & rate & delay  & even \\
\midrule
B1 flat             & \TODO{} & \TODO{} & \TODO{} & \TODO{} & \TODO{} \\
B2 mission deadline & \TODO{} & \TODO{} & \TODO{} & \TODO{} & \TODO{} \\
B3 bounded LTL      & \TODO{} & \TODO{} & \TODO{} & \TODO{} & \TODO{} \\
C phase ladder      & \TODO{} & \TODO{} & \TODO{} & \TODO{} & \TODO{} \\
\bottomrule
\end{tabular}
\end{table}

\TODO{E1, pick-and-place: Table~\ref{tab:main}; does the held-out
false-alarm rate respect $\alpha$; delay of C against B2 and B3;
localisation accuracy; behaviour on almost-failed episodes; one figure
of delay against false-alarm rate as $\alpha$ varies.}

\TODO{E2, navigation: the wedged-while-\texttt{following} case, in which
B4 and any planner-status proposition are silent and only bounds
grounded in the robot's own motion fire; delay per arm.}

\TODO{E3, cost: per-tick monitor time, split into automaton steps, guard
evaluation and proposition grounding.}

% =====================================================================
\section{Limitations and Conclusion}
\label{sec:conc}
% =====================================================================

\textbf{Limitations.} The main data are recordings, so the effect of
intervention on task success is not measured; break-even false-alarm
rate and execution saved are the proxies. Teleoperated runs have no
controller status, so independence is exercised only in navigation. The
calibration guarantee assumes exchangeable runs and is checked, not
proved, for new operators and scenes. Our shipped navigation
specification still grounded progress in the planner's status, which is
the very coupling this paper argues against and is why only its deadline
catches the motivating case; results use progress derived from the
robot's own motion. Missing propositions are not yet handled uniformly
(Section~\ref{sec:spec}). One robot, two skills, sequential phases only.

\textbf{Conclusion.} The formula that states a skill's goal can be
confirmed and never refuted, so the commonest failure of a robot skill
--- making no progress --- is invisible to a flat temporal
specification, and is left to the controller to notice about itself.
Declaring a calibrated bound per phase in the specification puts that
judgement in an independent monitor, fixes a horizon on its verdict, and
attaches a stage, a guard and a response to the fault. LTL is kept
alongside for the invariants and for the permanent verdict a supervisor
can latch. \TODO{one sentence of result.}

\bibliographystyle{IEEEtran}
\bibliography{refs}

\end{document}
